\section{Evidence Availability and Reporting Examples}\label{supp:reporting-examples}
\subsection{Eligibility and evidence-availability rules}
The descriptive tally uses the existing forty-four setting records from fifteen direct sources. It adds no retrieval or independent coding. The artifact's \path{transfer-synthesis/setting-availability.csv} carries each record's original access description, intervention, adaptation and utility evidence, source version and locator together with the new grouping. \path{availability-summary.json} gives the rules and exact membership behind every count.

Access is classified by the rights used in the evaluated setting: transit metadata (22 settings, 6 sources), addressed-recipient observations (10, 5), active intervention/probing (11, 6), or a downstream preference analysis (1, 1). These classes partition settings; source groups overlap. AgentPrint's persona analysis uses an inferred preference representation and is kept separate from a new packet-observer test. CHeaT's planted-token tests use the defender's active setup and knowledge. Combining access classes would grant an observer information from a different experiment.

The privacy denominator comprises \texttt{WRA-decoy}, \texttt{WRA-visibility}, \texttt{CG-transform} and \texttt{CG-live}: four query/workflow protection settings from two sources. Protected-condition fitting or an adaptive comparison is located for the two communication-graph settings. This records the inspected adaptation, not an optimal attacker. An original-task quality/goal measure for the same intervention is located for WRA visibility. Extra queries, attacker budgets, bytes and latency alone do not satisfy that field. Hence adaptation is located in 2/4, original-task utility in 1/4, and both in 0/4. Within the sole executed-cover setting, original-task utility is unlocated (0/1). WRA's utility does not transfer from the visibility condition to decoy execution.

Three identity-obfuscation records (AgentTrace's obfuscation setting and the two Fayolle fingerprinting settings with stealth configurations), two adversarial-Agent blocking records (Fayolle and CHeaT), and thirty-four records without a protection comparison are outside that denominator. The communication-graph composition setting re-measures the CG-transform properties on the wire-faithful corpus view and is counted with CG-transform. They remain in the corpus. ``Not located'' applies to retained evidence, while ``outside denominator'' states a different question; neither is a negative result. No field-completion fraction is used to score studies or validate the checklist. The table below states the four included decisions explicitly.

\begin{table*}[t]\centering\small
\caption{Setting-level basis for the query/workflow privacy denominator. Original record versions are retained; the official WRA text supplies the updated table locator.}\label{tab:availability-settings}
\begin{tabularx}{\textwidth}{@{}P{.18\textwidth}P{.21\textwidth}P{.24\textwidth}Y@{}}\toprule
Setting & Adaptation evidence & Original-task utility evidence & Source locator and boundary \\\midrule
WRA decoys & Matched protected ICL examples/refitting unresolved. & Not located for the decoy condition. & Official USENIX text, Sec.~8.1, Table~8; selected test-side mixing code. Five extra research traces do not supply a task-quality result. \\\addlinespace[3pt]
WRA visibility & Protected-condition refitting not located. & Located: report utility for the visibility intervention. & Sec.~8.2, Tables~10--11. Not a full live rerouting cost evaluation. \\\addlinespace[3pt]
Communication graphs, transformations & Located: protected fitting and adaptive comparison. & Not located. & 2606.07150v3, Secs.~9.3, 9.8 and 9.12 (composition corpus). Transformed captures do not measure a live task outcome. \\\addlinespace[3pt]
Communication graphs, live OHTTP & Located: fixed-family RF refit per protected capture. & Not located; bytes and latency are transmission costs. & 2606.07150v3, Sec.~9.7 and selected live-evaluation code. The broader attack-search output is a separate setting. \\\bottomrule
\end{tabularx}\end{table*}

\subsection{Checklist and completed examples}
Mark a box when its applicable fields and missing measurements are recorded.

\begin{list}{$\square$}{\setlength{\leftmargin}{1.4em}\setlength{\labelwidth}{1em}\setlength{\labelsep}{.4em}\setlength{\itemsep}{3pt}\setlength{\parsep}{0pt}}
\item \textbf{C1. Observer and links.} Name the observed interfaces, available metadata/content, access level and any supplied task boundaries.
\item \textbf{C2. Target and ground truth.} State whether labels denote intent, action, success, identity or risk; identify the policy reference for any authorization claim.
\item \textbf{C3. Population and failures.} Give denominators and the treatment of crashes, refusals, retries and incomplete executions; report the relevant execution outcome.
\item \textbf{C4. Split and transfer.} Name the split unit and each changed deployment axis; distinguish a supplied execution from an event found in a stream.
\item \textbf{C5. Attacks and adaptation.} State features, learners and still-visible links; compare unchanged attacks with attacks trained on protected executions.
\item \textbf{C6. Matched protection and cost.} Bind residual inference, original-task utility, added calls, bytes and delay to the same intervention and execution population.
\item \textbf{C7. Metrics and uncertainty.} Specify metric, denominator and uncertainty assumptions; for monitoring, report conditional error rates and the intended prevalence.
\item \textbf{C8. Source and availability.} Identify the source version and result locator, accessible artifacts and the scope of any code or data checks.
\end{list}


\begin{table*}[t]\centering\small
\caption{Relationship to existing evaluation guidance. These are field correspondences, not an empirical comparison of checklists. Arp P-numbers refer to Section~2 of the published paper~\cite{arp2022dos}.}\label{tab:guideline-crosswalk}
\setlength{\tabcolsep}{4pt}
\begin{tabularx}{\textwidth}{@{}P{.14\textwidth}P{.37\textwidth}Y@{}}\toprule
Entries & Established obligation & Agent-specific specification used here \\\midrule
C1, C4 & Realistic observer information and splits (Juarez et al.; Mathews et al.)~\cite{juarez2014critical,mathews2023critical}. & Separate transit, recipient and active access; identify supplied stage boundaries and association across model/tool links. \\\addlinespace[3pt]
C2, C3 & Sampling and label validity (Arp P1--P2). & Distinguish requested intent, logged action and checked outcome; state exclusions, null responses and recoverable errors. \\\addlinespace[3pt]
C4, C5 & Suitable comparisons (Arp P6), changed conditions and adaptive evaluation~\cite{juarez2014critical,mathews2023critical,tramer2020adaptive}. & Name composition, topology, scaffold or model changes; distinguish fixed algorithm from refitting on protected traffic. \\\addlinespace[3pt]
C6 & Protection overhead and realistic evaluation~\cite{dyer2012peekaboo,mathews2023critical}; deployment scope (Arp P9). & Bind original-task utility to executed cover, and bytes/latency to the same protection condition and population. \\\addlinespace[3pt]
C7 & Application-relevant metrics and prevalence (Arp P7--P8). & Report event detection separately from supplied-unit classification; attach uncertainty and base rates to the named monitoring target. \\\addlinespace[3pt]
C8 & Evidence record supporting the correspondences above. & Give source version, result locator and accessible artifacts; distinguish source inspection from experimental reproduction. \\
\bottomrule\end{tabularx}\end{table*}

These forms apply the Section~8 checklist to three selected settings already in the evidence matrix: WRA's five-decoy condition, the communication-graph live OHTTP condition and ANT's risk test. The authors filled them from the recorded source passages and code checks. They introduce no new source experiment or independent analyst evaluation.

\paragraph{Reading a filled form} \emph{Reported} identifies supplied evidence; \emph{partial} identifies a supplied component and its specific gap. \emph{Not located} applies to the inspected passages or mapped record, not all possible source artifacts. \emph{Not applicable} means that the field is outside this setting's claim. A checkmark in a new study's blank checklist means the field has been filled or explicitly dispositioned; it is not a pass mark. The examples distinguish an unavailable interval in the current extraction from a claim that the source never reports uncertainty.

\begin{table*}[t]\centering\small
\caption{Completed forms for all eight reporting fields. Versions and source locations are given in C8; entries preserve the source-specific observer and denominator.}\label{tab:filled-checklists}
\setlength{\tabcolsep}{4pt}
\begin{tabularx}{\textwidth}{@{}P{.11\textwidth}YYY@{}}\toprule
Field & WRA: decoy recovery & Communication graphs: live OHTTP & ANT: benchmark risk \\\midrule
C1 & Reported: research-task domain mixture; passive observation and supplied task boundary. & Reported: live link capture with relay/trust assumptions; supplied workflow. & Reported: passive packet metadata with supplied episode boundaries. \\ \addlinespace[5pt]
C2 & Reported: original query meaning is the target. Application authorization is not applicable. & Reported: one of six workflow labels. Application authorization is not applicable. & Partial: benign/misuse task labels are explicit; per-execution principal/approval ground truth is not located. \\ \addlinespace[5pt]
C3 & Partial: real-plus-decoy research is specified; matched original-task outcome and execution-failure handling are not located. & Partial: the live population is distinct; matched task-goal outcome and failure selection are not located. & Partial: 20 positive and 189 negative test episodes; refusal-only Attack Code Generation is excluded, without general failure stratification. \\ \addlinespace[5pt]
C4 & Partial: a supplied-task decoy comparison; the mapped record does not supply a complete split-unit description or a deployment-transfer test. & Reported: fold-wise classification within each live protected capture; this six-class population is distinct from the nine-class composition test. & Reported: within-task episode split and 209 test episodes; this is not a new-task or authorization-transfer test. \\ \addlinespace[5pt]
C5 & Partial: test-side decoy mixing is checked; matched protected ICL examples/refitting remain unresolved. & Partial: fixed-family RF refitting is checked; the broader attack-search output is not bound to this live setting. & Not applicable: protected-condition adaptation, because this risk setting has no privacy intervention. Thirteen unprotected classifiers are reported. \\ \addlinespace[5pt]
C6 & Partial: five extra research traces and recovery .770 to .565. Matched decoy bytes, delay and original-task utility are not located. & Partial: accuracy .30 and about 13-fold bytes/11-fold latency share the live condition; matched task utility is not located. & Not applicable: a protection/cost comparison. The setting evaluates benchmark-risk recognition. \\ \addlinespace[5pt]
C7 & Partial: functional recovery is recorded; a matched uncertainty interval is not transcribed in this setting record. Deployment alarm prevalence is not this reconstruction metric. & Partial: six-class accuracy and cost ratios are recorded; a matched uncertainty interval is not transcribed in this setting record. & Reported: class totals and rounded accuracy, F1, FPR and recall. This SoK adds conditional binomial reference bands; raw clustering and deployment shift remain unknown. \\ \addlinespace[5pt]
C8 & Reported: official USENIX 2026 PDF, Sec. 8.1, decoy Table 8; visibility Table 11 is separate. Selected implementation checks are in Section~\ref{supp:sources}. & Reported: arXiv 2606.07150v3, Sec. 9.7 and the selected live evaluation path in Section~\ref{supp:sources}; no independent rerun. & Reported: arXiv 2610.06514v1, Secs. 3.1, 4.1--4.2, App. A.1, Table 3. Public-table reconstruction is in the artifact\textquotesingle s \texttt{impact-analysis/} directory. \\ \addlinespace[5pt]
\bottomrule\end{tabularx}\end{table*}

The completed examples retain different open fields. WRA needs matched decoy adaptation and utility before its recovery reduction supports a complete privacy--utility point. The live OHTTP example already supports fixed-family protected-data refitting and transmission cost, while leaving task utility unresolved. ANT supports a benchmark-risk classifier comparison; an authorization judgment would additionally need the relevant execution policy and principal. None of these distinctions requires changing a reported score.
